\documentclass[11pt,a4paper]{article}
\usepackage[margin=1in]{geometry}
\usepackage{amsmath,amssymb,amsthm}
\usepackage{algorithm}
\usepackage{algpseudocode}
\usepackage{booktabs}
\usepackage{hyperref}

\theoremstyle{definition}
\newtheorem{definition}{\textbf{Definition}}[section]
\newtheorem{theorem}{\textbf{Theorem}}[section]
\newtheorem{lemma}[theorem]{\textbf{Lemma}}
\newtheorem{remark}{\textbf{Remark}}[section]

\title{\textbf{Semi-Structured 2:4 Fine-Grained Sparsity for Hardware Sparse Tensor Cores}}
\author{\textbf{Team Scaibu} \\ \small Email: \texttt{jaydeep.9664920749@gmail.com} \\ \small \textbf{Architecture and Core Engine Team}}
\date{\textbf{October 2026}}

\begin{document}
\maketitle

\section{Definitions and Cognitive Mental Models}

\subsection{Formal Mathematical Definition}
\textbf{Definition (Semi-Structured 2:4 Fine-Grained Sparsity):}
While unstructured sparsity provides theoretical model compression, modern graphics processing units (GPUs) require regular, aligned data access to avoid memory divergence and load imbalances. Beginning with the NVIDIA Ampere (A100) architecture and continuing through Hopper (H100) and Blackwell (B200), hardware Sparse Tensor Cores natively accelerate \textbf{2:4 Fine-Grained Semi-Structured Sparsity} (Mishra et al., 2021).

Let $\mathbf{v} \in \mathbb{R}^N$ be a flattened dense parameter vector where $N$ is a multiple of 4 ($N = 4B$, $B \in \mathbb{N}$). The vector is partitioned into $B$ disjoint, contiguous four-element blocks:
{\boldmath
\begin{equation}
\mathbf{v} = \left[ \mathbf{b}_1, \mathbf{b}_2, \dots, \mathbf{b}_B \right], \quad \mathbf{b}_k = \left( v_{4k-3}, \, v_{4k-2}, \, v_{4k-1}, \, v_{4k} \right)^\top \in \mathbb{R}^4
\end{equation}
}
Under the 2:4 sparsity constraint, exactly \textbf{two out of four} elements in every block $\mathbf{b}_k$ must be zero:
{\boldmath
\begin{equation}
\|\mathbf{b}_k\|_0 = 2, \quad \forall k \in \{1, \dots, B\}
\end{equation}
}
The projection operator $\Pi_{2:4}: \mathbb{R}^4 \to \mathbb{R}^4$ selects the two coordinates with largest absolute magnitudes $|v_{4k-3+\pi(1)}| \ge |v_{4k-3+\pi(2)}| \ge \dots$:
{\boldmath
\begin{equation}
\Pi_{2:4}(\mathbf{b}_k)_j = \begin{cases} b_{k, j} & \text{if } j \in \{ \pi(1), \pi(2) \} \\ 0 & \text{if } j \in \{ \pi(3), \pi(4) \} \end{cases}
\end{equation}
}
The resulting sparse tensor achieves an exact $50\%$ memory reduction by packing the two non-zero floating-point values into half-sized memory buffers accompanied by a 2-bit offset index for hardware systolic decompression, doubling GEMM throughput.

\subsection{Expanded Non-Technical Definition (Intuition for Anyone)}
\begin{enumerate}
    \item \textbf{The Buddy System for GPU Chips}: Normal random pruning is like letting people sit randomly in an auditorium; the ushers waste enormous time finding empty chairs. 2:4 sparsity creates a strict rule: "Every group of 4 seats must contain exactly 2 empty chairs and 2 occupied chairs."
    \item \textbf{Instant 2$\times$ Speed Boost on Real Hardware}: Because GPU silicon is physically built to expect this 2-out-of-4 pattern, the hardware tensor cores skip the zeros automatically and multiply twice as fast, cutting inference time in half with zero memory latency penalty.
    \item \textbf{Minimal Accuracy Drop}: Unlike structured channel pruning (which deletes entire attention heads or hidden layers, causing severe accuracy drops), 2:4 sparsity distributes the pruning evenly across every micro-block in the network, preserving the representational rank of every layer.
\end{enumerate}

\subsection{Child-Friendly Mental Model (Physical Metaphor)}
Imagine packing lunchboxes:
\begin{itemize}
    \item \textbf{Messy Packing}: Randomly tossing snacks in some boxes and leaving others overflowing.
    \item \textbf{2:4 Standard}: Every box has 4 slots, and you always put exactly 2 snacks and leave 2 slots empty. The conveyor belt moves at double speed because every box is identical and balanced!
\end{itemize}

\section{Core Mathematical Equations and Definitions}

\subsection{Hardware 2-Bit Index Combinatorial Encoding}
{\boldmath
\begin{equation}
\mathcal{C} = \binom{4}{2} = \frac{4!}{2! \, 2!} = 6 \le 2^3
\end{equation}
}
\begin{itemize}
    \item \textbf{Formal Definition}: The cardinality of the configuration space of selecting 2 active indices out of 4 contiguous offsets.
    \item \textbf{What It Means}: There are exactly 6 possible index pairs: $(0,1), (0,2), (0,3), (1,2), (1,3), (2,3)$.
    \item \textbf{Why It Is Important}: Hardware encodes the pattern in a tiny 2-bit or 4-bit metadata nibble, adding negligible storage overhead ($< 12.5\%$) compared to full FP16/FP8 weights.
\end{itemize}

\subsection{Block Frobenius Projection Error}
{\boldmath
\begin{equation}
\|\mathbf{b}_k - \Pi_{2:4}(\mathbf{b}_k)\|_2^2 = |v_{4k-3+\pi(3)}|^2 + |v_{4k-3+\pi(4)}|^2
\end{equation}
}
\begin{itemize}
    \item \textbf{Formal Definition}: Exact Euclidean squared reconstruction error incurred on block $k$ under the magnitude-optimal 2:4 projection.
    \item \textbf{What It Means}: Sum of the squared magnitudes of the two smallest elements in the 4-element window.
    \item \textbf{Why It Is Important}: Proves that sorting by absolute value minimizes the total Frobenius reconstruction error across the entire tensor.
\end{itemize}

\section{Rigorous Mathematical Analysis Checklist}

\begin{itemize}
    \item \textbf{Notation}: $N = 4B$ (vector length), $B$ (block count), $\mathbf{b}_k \in \mathbb{R}^4$ (4-element window), $\Pi_{2:4}$ (projection operator).
    \item \textbf{Epistemic Status}: Hardware standard in NVIDIA TensorRT, cuSPARSELt, and PyTorch Sparse Semi-Structured.
    \item \textbf{Dependency Graph}: Dense Vector $\to$ 4-Element Block Chunking $\to$ Absolute Value Ranking $\to$ Top-2 Selection $\to$ Compressed Values \& Index Manifest.
    \item \textbf{Dimensions}: Input vector $\mathbb{R}^N$; compressed values $\mathbb{R}^{N/2}$; index list $\mathbb{N}_0^{B \times 2}$.
    \item \textbf{Regimes}: Production deployment on NVIDIA A100, H100, L40S, and B200 Tensor Core architectures.
    \item \textbf{Invariants}: Exact 50\% compression ratio: $\|\Pi_{2:4}(\mathbf{v})\|_0 \equiv \frac{N}{2}$ for all valid inputs $N \equiv 0 \pmod 4$.
    \item \textbf{Isomorphisms}: Isomorphic to hard block-thresholding pursuit in compressive sensing with structured dictionary constraints.
\end{itemize}

\section{Theoretical Theorems and Formal Proofs}

\begin{theorem}[Optimality of Greedy Magnitude Selection Under 2:4 Frobenius Projection]
Let $\mathbf{b} = (v_1, v_2, v_3, v_4)^\top \in \mathbb{R}^4$. Let $\mathcal{S}_{2:4} = \{ \mathbf{u} \in \mathbb{R}^4 : \|\mathbf{u}\|_0 \le 2 \}$ denote the set of all vectors in $\mathbb{R}^4$ with at most two non-zero coordinates. The Euclidean projection $\Pi_{2:4}(\mathbf{b}) = \operatorname{argmin}_{\mathbf{u} \in \mathcal{S}_{2:4}} \|\mathbf{b} - \mathbf{u}\|_2^2$ is uniquely achieved by retaining the two coordinates of largest absolute magnitude:
{\boldmath
\begin{equation}
\min_{\mathbf{u} \in \mathcal{S}_{2:4}} \|\mathbf{b} - \mathbf{u}\|_2^2 = |v_{\pi(3)}|^2 + |v_{\pi(4)}|^2
\end{equation}
}
where $|v_{\pi(1)}| \ge |v_{\pi(2)}| \ge |v_{\pi(3)}| \ge |v_{\pi(4)}|$.
\end{theorem}

\begin{proof}
For any $\mathbf{u} \in \mathcal{S}_{2:4}$, let $\operatorname{supp}(\mathbf{u}) = \{j \in \{1, 2, 3, 4\} : u_j \neq 0\}$ with $|\operatorname{supp}(\mathbf{u})| \le 2$.
The squared distance is:
\begin{equation}
\|\mathbf{b} - \mathbf{u}\|_2^2 = \sum_{j \in \operatorname{supp}(\mathbf{u})} (v_j - u_j)^2 + \sum_{j \notin \operatorname{supp}(\mathbf{u})} v_j^2
\end{equation}
To minimize this expression over $\mathbf{u}$, for any fixed choice of $\operatorname{supp}(\mathbf{u})$, the optimal choice of non-zero entries is clearly $u_j = v_j$ for all $j \in \operatorname{supp}(\mathbf{u})$, which zeroes the first sum:
\begin{equation}
\min_{\mathbf{u} \text{ with supp}(\mathbf{u})} \|\mathbf{b} - \mathbf{u}\|_2^2 = \sum_{j \notin \operatorname{supp}(\mathbf{u})} v_j^2
\end{equation}
Thus, the problem reduces to choosing a subset of at most two coordinates to retain (or equivalently, choosing a subset of two coordinates $\mathcal{I}_{\text{zero}}$ to zero) such that the sum of squares of zeroed coordinates is minimized:
\begin{equation}
\min_{\mathcal{I}_{\text{zero}} \subset \{1, 2, 3, 4\}, \, |\mathcal{I}_{\text{zero}}| = 2} \sum_{j \in \mathcal{I}_{\text{zero}}} v_j^2
\end{equation}
The sum of two elements chosen from $\{v_1^2, v_2^2, v_3^2, v_4^2\}$ is strictly minimized by selecting the two smallest values. Since $v_j^2 = |v_j|^2$, sorting by absolute value yields $|v_{\pi(3)}|^2 + |v_{\pi(4)}|^2$, proving that keeping the top-2 absolute elements achieves the global minimum projection error.
\end{proof}

\begin{theorem}[Deterministic Linear Compression and Storage Bound]
Let $\mathbf{V} \in \mathbb{R}^{M \times K}$ be a weight matrix where $K \equiv 0 \pmod 4$. The 2:4 structured representation reduces weight storage volume by exactly 50\%, requiring total bit consumption:
{\boldmath
\begin{equation}
\mathcal{B}_{2:4} = M \cdot K \cdot \left( \frac{b_{\text{val}}}{2} + \frac{b_{\text{idx}}}{4} \right)
\end{equation}
}
where $b_{\text{val}}$ is the bit-width of weight precision (e.g., 16 bits for FP16, 8 bits for FP8) and $b_{\text{idx}} = 4$ bits per block metadata nibble. For FP16 ($b_{\text{val}}=16$), the effective compression ratio is $56.25\%$, and for FP8 ($b_{\text{val}}=8$), effective storage is $62.5\%$ of dense footprint while accelerating matrix multiplication by $2\times$.
\end{theorem}

\begin{proof}
Let the total number of parameter elements be $P = M \cdot K$.
Under 2:4 sparsity, the matrix is partitioned into $B = P / 4$ blocks.
Each block contains exactly 2 surviving values and 1 metadata nibble encoding the index pair:
\begin{equation}
\mathcal{B}_{\text{block}} = 2 \cdot b_{\text{val}} + b_{\text{idx}}
\end{equation}
Multiplying across all $B = P / 4$ blocks:
\begin{equation}
\mathcal{B}_{2:4} = \frac{P}{4} \left( 2 b_{\text{val}} + b_{\text{idx}} \right) = P \cdot \left( \frac{b_{\text{val}}}{2} + \frac{b_{\text{idx}}}{4} \right)
\end{equation}
Dense storage consumes $\mathcal{B}_{\text{dense}} = P \cdot b_{\text{val}}$ bits.
The compression fraction is:
\begin{equation}
\rho = \frac{\mathcal{B}_{2:4}}{\mathcal{B}_{\text{dense}}} = \frac{\frac{b_{\text{val}}}{2} + \frac{b_{\text{idx}}}{4}}{b_{\text{val}}} = \frac{1}{2} + \frac{b_{\text{idx}}}{4 b_{\text{val}}}
\end{equation}
For FP16 ($b_{\text{val}}=16, b_{\text{idx}}=4$): $\rho = \frac{1}{2} + \frac{4}{64} = 0.5 + 0.0625 = 0.5625$ (56.25\%).
For FP8 ($b_{\text{val}}=8, b_{\text{idx}}=4$): $\rho = \frac{1}{2} + \frac{4}{32} = 0.5 + 0.125 = 0.625$ (62.5\%).
Hardware Sparse Tensor Cores bypass reading the zeroes entirely from DRAM/SRAM, achieving an exact $2\times$ theoretical arithmetic acceleration in the systolic multiplier pipeline.
\end{proof}

\section{Architectural Questions and Real-World Trade-Offs}

\subsection{Architectural Question 1: Sparse Tensor Core Hardware Mechanics}
\textbf{Question:} How does the physical hardware in NVIDIA Ampere/Hopper execute 2:4 Sparse Tensor Multiplications?\\
\textbf{Answer:} In dense GEMM, a $16 \times 16$ tile processes 256 multiply-accumulate (MAC) operations per clock cycle. With 2:4 Sparse Tensor Cores, the matrix memory controller reads the 2 non-zero weights per block along with the 2-bit metadata index. A hardware multiplexer inside the tensor core routes the 2 sparse weights to match their corresponding activation vectors, multiplying a compressed $16 \times 8$ weight matrix with a $8 \times 16$ activation tile in half the time, doubling peak TFLOPS.

\subsection{Architectural Question 2: Retraining with 2:4 Sparsity (SR-STE)}
\textbf{Question:} Can foundation models be converted to 2:4 sparsity post-training, or do they require retraining?\\
\textbf{Answer:} Pruning a dense pre-trained model directly to 2:4 sparsity in one shot incurs a small perplexity penalty ($\sim 0.3 - 0.8$ perplexity on WikiText). To completely recover dense accuracy, Sparse-Refined Straight-Through Estimator (SR-STE) fine-tuning is applied for 1--2\% of the original training budget. During forward passes, weights are projected onto the 2:4 mask; during backward passes, gradients update the dense unmasked weights, allowing parameters to fluidly swap into and out of the active set.

\section{Algorithmic Implementation Specification}

\begin{algorithm}[H]
\caption{Semi-Structured 2:4 Fine-Grained Sparsity Projection}
\begin{algorithmic}[1]
\Require Dense vector $\mathbf{v} \in \mathbb{R}^N$ ($N \equiv 0 \pmod 4$), importance metric $\mathcal{M} \in \{\text{"magnitude"}, \text{"signed"}\}$
\Ensure Sparse vector $\mathbf{v}_{\text{sparse}} \in \mathbb{R}^N$, compressed values $\mathbf{c} \in \mathbb{R}^{N/2}$, index pairs $\mathcal{I} \in \mathbb{N}_0^{(N/4) \times 2}$, block count $B$
\State $N \gets \operatorname{len}(\mathbf{v})$
\State $B \gets N / 4$
\State Initialize $\mathbf{v}_{\text{sparse}} \gets [0.0]^N, \quad \mathbf{c} \gets [\,], \quad \mathcal{I} \gets [\,]$
\For{$b = 0$ \textbf{to} $B - 1$}
    \State $\text{base} \gets b \cdot 4$
    \State $\mathbf{g} \gets [ \mathbf{v}[\text{base}], \, \mathbf{v}[\text{base}+1], \, \mathbf{v}[\text{base}+2], \, \mathbf{v}[\text{base}+3] ]$
    \If{$\mathcal{M} = \text{"signed"}$}
        \State $\boldsymbol{\pi} \gets \operatorname{argsort\_descending}([ g_0, g_1, g_2, g_3 ])$
    \Else
        \State $\boldsymbol{\pi} \gets \operatorname{argsort\_descending}([ |g_0|, |g_1|, |g_2|, |g_3| ])$
    \EndIf
    \State $\text{top\_two} \gets \operatorname{sort\_ascending}(\boldsymbol{\pi}[0 : 2])$
    \State Append $\text{top\_two}$ to $\mathcal{I}$
    \For{\textbf{each} $\text{off} \in \text{top\_two}$}
        \State $\text{val} \gets \mathbf{g}[\text{off}]$
        \State $\mathbf{v}_{\text{sparse}}[\text{base} + \text{off}] \gets \text{val}$
        \State Append $\text{val}$ to $\mathbf{c}$
    \EndFor
\EndFor
\State \Return $(\mathbf{v}_{\text{sparse}}, \mathbf{c}, \mathcal{I}, 0.5, B)$
\end{algorithmic}
\end{algorithm}

\section{Big-$\Theta$ Complexity Analysis}

\begin{itemize}
    \item \textbf{4-Element Block Iteration}: $\Theta(N/4) = \Theta(N)$ window steps.
    \item \textbf{Local Offset Ranking}: $\Theta(1)$ constant time sorting (4 elements per block).
    \item \textbf{Compressed Array Extraction}: $\Theta(N/2) = \Theta(N)$ linear memory write passes.
    \item \textbf{Total Time Complexity}: $\Theta(N)$ strictly linear with respect to parameter count.
    \item \textbf{Space Complexity}: $\Theta(N)$ output storage for sparse and compressed representations.
\end{itemize}

\section{Single Canonical Repository Reference}

The complete deterministic scalar implementation, verification suite, and unit test benchmarks are published at:
\begin{center}
\url{https://github.com/Chief-Strategist-J/llm-obs-infra}
\end{center}

\end{document}
